% Owner: Abhishek Kumar Patel. Proposed Solution: single-cycle baseline.
\subsection{Single-Cycle Baseline}
To judge what pipelining gains, the pipelined processor is compared with a
single-cycle processor that implements the same instruction set. This follows
the usual design order: first a correct single-cycle datapath, then its
partition into pipeline stages~\cite{pattersonRV,harrisRV}. The baseline is the
module \texttt{riscv\_single\_cycle\_top}. It has the same 29-instruction subset,
the same separate instruction and data memories (1024 words each, combinational
read, synchronous write) and the same instance names as the pipeline
(\texttt{u\_imem}, \texttt{u\_reg\_file}, \texttt{u\_dmem}), so the same testbenches
can read both designs.

\subsubsection{Datapath}
Every instruction is fetched, decoded, executed, accesses memory and writes its
result back in one clock cycle (Fig.~\ref{fig:single}). The baseline reuses the
pipeline modules unchanged: \texttt{pc\_reg} (always enabled),
\texttt{instr\_mem}, \texttt{control\_unit}, \texttt{imm\_gen},
\texttt{alu\_control}, \texttt{alu}, \texttt{branch\_unit} and \texttt{data\_mem}.
Only the glue logic differs:
\begin{itemize}
  \item ALU operand A is \texttt{rd1}, or zero for LUI so that the ALU computes
        $0 + imm$; operand B is the immediate or \texttt{rd2}, selected by
        \texttt{alu\_src}.
  \item The next PC is \texttt{branch\_target} when \texttt{pc\_src} is
        asserted and $pc + 4$ otherwise. The branch unit compares the register
        values read in the same cycle, so a taken branch or jump costs no
        extra cycles and nothing has to be flushed.
  \item The value written back is the loaded word for LW, the link address
        $pc + 4$ for JAL and JALR, and the ALU result otherwise.
\end{itemize}
There are no pipeline registers, no forwarding unit and no hazard detection
unit. With one instruction in flight, no data, control or structural hazard can
occur.

\begin{figure*}[!t]
\centering
\begin{tikzpicture}[
  font=\scriptsize, >={Stealth[length=1.5mm]},
  blk/.style={draw, rounded corners=2pt, minimum width=1.5cm, minimum height=0.75cm, align=center, fill=white},
  mux/.style={draw, fill=white, trapezium, trapezium angle=72, shape border rotate=270, minimum height=0.95cm, inner sep=1pt},
  w/.style={->, thick},
  c/.style={->, dashed, black!60},
  lbl/.style={font=\tiny, inner sep=1pt}
]
% blocks
\node[mux] (pcm) at (-0.4,0) {};
\node[blk, minimum width=0.8cm] (pc) at (0.7,0) {PC};
\node[blk] (im) at (2.55,0) {Instr.\\Memory};
\node[blk] (cu) at (6.0,1.9) {Control\\Unit};
\node[blk, minimum height=1.15cm] (rf) at (6.0,0) {Register\\File};
\node[blk] (ig) at (6.0,-1.9) {Imm.\\Generator};
\node[mux] (mb) at (8.35,-0.55) {};
\node[blk, minimum width=1.0cm, minimum height=1.5cm] (alu) at (9.85,-0.1) {ALU};
\node[blk, minimum width=1.5cm] (bu) at (9.85,2.15) {Branch\\Unit};
\node[blk, minimum height=1.15cm] (dm) at (12.4,-0.1) {Data\\Memory};
\node[mux] (wb) at (14.5,-0.1) {};
\node[blk, minimum width=0.9cm, minimum height=0.5cm] (p4) at (2.55,1.9) {+4};
% fetch
\draw[w] (pcm.east) -- (pc.west);
\draw[w] (pc.east) -- (im.west);
\draw[w] (pc.east) ++(0.15,0) |- (p4.west);
\fill (pc.east) ++(0.15,0) circle (1.1pt);
\draw[w] (p4.east) -- (3.4,1.9) -- (3.4,2.85) -- (-0.75,2.85) -- (-0.75,0.25) -- ($(pcm.west)+(0,0.25)$);
\node[lbl] at (1.6,3.0) {pc + 4};
% instruction bus
\draw[w] (im.east) -- (4.5,0);
\draw[w] (4.5,1.9) -- (cu.west);
\draw[w] (4.5,0) -- (rf.west);
\draw[w] (4.5,-1.9) -- (ig.west);
\draw (4.5,1.9) -- (4.5,-1.9);
\fill (4.5,0) circle (1.1pt);
\node[lbl] at (3.75,0.2) {instr};
% register file to ALU
\draw[w] ($(rf.east)+(0,0.3)$) -- node[lbl,above]{rd1} ($(alu.west)+(0,0.4)$);
\draw[w] ($(rf.east)+(0,-0.3)$) -- node[lbl,above,pos=0.35]{rd2} (mb.west |- 0,-0.3);
\draw[w] (ig.east) -- node[lbl,above,pos=0.3]{imm} (mb.west |- 0,-1.9) -- (mb.west |- 0,-0.8);
\draw[w] (mb.east) -- ($(alu.west)+(0,-0.45)$);
% ALU outputs
\draw[w] (alu.east) -- node[lbl,above]{result} (dm.west);
\draw[w] (rf.south east) ++(0.35,0) -- ++(0,-0.95) -- (11.3,-1.9) -- (11.3,-0.55) -- ($(dm.west)+(0,-0.45)$);
\node[lbl] at (10.2,-2.05) {rd2 (store data)};
\draw[w] (dm.east) -- node[lbl,above]{rdata} ($(wb.west)+(0,0.3)$);
\draw[w] (11.0,-0.1) -- ++(0,-1.25) -- (13.45,-1.35) -- (13.45,-0.4) -- ($(wb.west)+(0,-0.3)$);
\fill (11.0,-0.1) circle (1.1pt);
% link address to write-back
\draw[w] (3.4,2.85) -- (14.1,2.85) -- (14.1,0.4) -- ($(wb.west)+(0,0.0)$) ;
\node[lbl] at (8.2,3.0) {pc + 4 (link address for JAL, JALR)};
% write-back to register file
\draw[w] (wb.east) -- (15.35,-0.1) -- (15.35,-2.65) -- (5.2,-2.65) -- (5.2,-0.45) -- ($(rf.west)+(0,-0.45)$);
\node[lbl] at (10.2,-2.8) {write-back data to register file};
% branch unit
\fill (7.9,0.3) circle (1.1pt);
\draw[w] (7.9,0.3) |- (bu.west);
\node[lbl,anchor=east,align=right] at (7.85,1.25) {rd1, rd2,\\pc, imm};
\draw[w] (bu.north) -- (9.85,3.35) -- (-1.05,3.35) -- (-1.05,-0.25) -- ($(pcm.west)+(0,-0.25)$);
\node[lbl] at (5.0,3.5) {branch\_target, pc\_src};
% control signals
\draw[c] (cu.east) -- ++(0.7,0);
\node[lbl,anchor=west] at (7.55,2.05) {control};
\end{tikzpicture}
\caption{Datapath of the single-cycle baseline. Every instruction passes through
all blocks in one clock cycle. There are no pipeline registers, forwarding paths
or hazard logic. The dashed line stands for the control signals that go from
the control unit to the ALU, memory and write-back multiplexers.}
\label{fig:single}
\end{figure*}

\subsubsection{Register File Without Bypass}
The pipeline's \texttt{reg\_file} returns the write-back value when the register
being written is read in the same cycle (write-first bypass), which removes the
WB-to-ID hazard. In a single-cycle design this bypass is harmful: the value to
write depends on the value being read, so an instruction such as
\texttt{add x8, x8, x9} creates a combinational loop from the read port through
the ALU and write-back back to the read port. This was found when the
simulation of the array-sum program never finished. The baseline therefore uses
\texttt{reg\_file\_nobypass}, the same $32\times32$ array and ports without the
bypass: a read returns the old value during the cycle and the new value is
stored on the clock edge.

\subsubsection{Timing and Cycle Count}
Because each instruction takes exactly one cycle, the cycle count equals the
instruction count, $C = N$, and CPI is 1 by construction. The price is the
clock period: it must cover the whole path from the program counter through the
instruction memory, register file, ALU and data memory back to the register
file write. The pipeline splits this path into five shorter stages. The
resulting clock frequencies are measured in the synthesis results, and both
effects are combined in the performance comparison.

\subsubsection{Verification of the Baseline}
The baseline runs the same six programs as the pipeline, with the same expected
register and memory values, so both designs must end in the same architectural
state. All 36 system checks and the 10 checks of the register-file unit test
pass. Table~\ref{tab:singlecheck} shows that the baseline executes exactly the
same number of instructions as the pipeline on every program.

\begin{table}[!t]
\centering
\caption{Single-Cycle Baseline on the Phase I Programs (Icarus Verilog)}
\label{tab:singlecheck}
\setlength{\tabcolsep}{5pt}
\renewcommand{\arraystretch}{1.08}
\begin{tabular}{@{}lccc@{}}
\toprule
\textbf{Program} & \textbf{Checks} & \textbf{$N$} & \textbf{Cycles} \\ \midrule
1 arithmetic and logic & 18 & 21 & 21 \\
2 load / store         & 4  & 7  & 7 \\
3 forwarding           & 3  & 6  & 6 \\
4 load-use             & 3  & 5  & 5 \\
5 branches and jumps   & 6  & 11 & 11 \\
6 array sum loop       & 2  & 40 & 40 \\ \bottomrule
\end{tabular}
\end{table}
